%
% Beweis des Prinzips von Liouville in Spezialfällen
%
\subsection{Beweis des Prinzips von Liouville in Spezialfällen
\label{buch:intalg:liouville:subsection:spezialfaelle}}
Das Prinzip von Liouville geht davon aus, dass eine Funktion $f$
in einem Funktionenkörper eine Stammfunktion $g$ in einer geeigneten
Erweiterung hat.
Es zeigt dann, dass diese Erweiterung nur Logarithmen von Elementen
in $F$ braucht und dass diese Logarithmen in der Stammfunktion nur 
linear vorkommen.
Beim Ableiten müssen die Terme, die Elemente der Erweiterung enthalten,
wieder verschwinden.
Man muss also für exponentielle und algebraische Erweiterungen zeigen,
dass diese beim Ableiten nicht verschwinden können.
Und für logarithmische Erweiterungen muss man zeigen, dass diese genau
dann verschwinden können, wenn sie linear in der Stammfunktion auftreten.

In diesem Abschnitt beweisen wir dies für eine Erweiterung, die von einem
einzigen zusätzlichen Element erzeugt wird.
Dies ist zwar nicht ausreichend für das Integral
\[
\int \frac{dx}{x^2+1}
=
\frac{1}{2i}\log(x+i)
-
\frac{1}{2i}\log(x-i)
=
\arctan x,
\]
welches zwei unabhängige logarithmische Erweiterungen braucht, aber
die Methode, die wir dabei kennenlernen werden, wird in
Abschnitt~\ref{buch:intalg:liouville:subsection:allgemein}
verallgemeinert, so dass sie auch für kompliziertere Integrale
gilt und damit das liouvillesche Prinzip bestätigt.

%
% Ein Ausdruck für die Stammfunktion
%
\subsubsection{Ein Ausdruck für die Stammfunktion}
In diesem Abschnitt ist $F$ ein Funktionenkörper mit Konstantenkörper
$K$ und $f\in F$ eine Funktion.
Es wird angenommen, dass es eine Stammfunktion $g$ von $f$ in einem
Erweiterungskörper $F(\vartheta)$ gibt.
$g\in F(\vartheta)$ bedeutet aber, dass $g$ ein rationaler Ausdruck
in $\vartheta$ ist.
Insbesondere kann $g$ als
\begin{equation}
g
=
a_0(\vartheta)
+
\frac{a(\vartheta)}{b(\vartheta)}
\label{buch:intalg:liouville:spezial:eqn:bruch}
\end{equation}
geschrieben werden, wobei $a_0(\vartheta)$, $a(\vartheta)$ und $b(\vartheta)$
Polynome mit Koeffizienten in $F$ sind.
Ausserdem darf angenommen werden, dass $a(\vartheta)$ und $b(\vartheta)$
teilerfremd sind, dass also der Bruch in
\eqref{buch:intalg:liouville:spezial:eqn:bruch}
vollständig gekürzt ist und dass
$\deg a(\vartheta)<\deg b(\vartheta)$ gilt.

Wir nehmen ausserdem an, dass $F$ und $F(\vartheta)$ den gleichen
Konstantenkörper haben.
Dies kann immer erreicht werden, indem man alle für das Integral
nötigen Konstanten bereits dem Körper $F$ hinzufügt, bevor man beginnt,
um nichtkonstante Funktionen zu erweitern.

Wir müssen zeigen, dass $\vartheta$ nur ein Logarithmus sein kann und dass
$\vartheta$ nur linear in $g$ vorkommen kann.

%
% Der Fall einer logarithmischen Erweiterung
%
\subsubsection{Der Fall einer logarithmischen Erweiterung}
Das Prinzip von Liouville besagt nicht nur, dass in der Stammfunktion nur
logarithmische Erweiterungen vorkommt, sondern dass sie sogar nur
linear mit konstanten Koeffizienten vorkommen.
In diesem Abschnitt betrachten wir eine einzelne logarithmische
Erweiterungen $F(\vartheta)\supset F$.
Das Prinzip von Liouville besagt in diesem Spezialfall, dass die
Stammfunktion von $f\in F$, falls sie in $F(\vartheta)$ liegt,
die Form $c_0+c_1\vartheta$ haben muss, wobei $c_1$ eine Konstante ist.

\begin{satz}
\label{buch:intalg:liouville:satz:logsimple}
Sei $F(\vartheta)$ eine logarithmische Erweiterung eines Funktionenkörpers
$F$.
$F$ und $F(\vartheta)$ haben den gleichen Konstantenkörper.
Weiter sei $f\in F$ und $g \in F(\vartheta)$ eine Stammfunktion.
Dann ist $g=c_0+c_1\vartheta$ mit $c_0\in F$ und $c_1\in K$.
\end{satz}

\begin{proof}
Der Nenner $b(\vartheta)$ des Ausdrucks 
\eqref{buch:intalg:liouville:spezial:eqn:bruch}
für $g$ kann in die irreduziblen Faktoren
$b_1(\vartheta),\ldots,b_m(\vartheta)$
zerlegt werden.
Weiter dürfen wir annehmen, dass die Faktoren $b_k(\vartheta)$
normierte Polynome sind.
Die zugehörige Partialbruchzerlegung hat daher die Form
\begin{align}
g
&=
a_0(\vartheta)
+
\sum_{k=1}^m
\sum_{i=1}^{r_k}
\frac{a_{ik}(\vartheta)}{b_k(\vartheta)^i},
\label{buch:intalg:liouville:speziell:eqn:f}
\intertext{wobei $a_{ik}(\vartheta)\in F(\vartheta)$ Polynome mit
$\deg a_{ik}(\vartheta) < \deg b_k(\vartheta)$ sind.
Die Ableitung von $g$ ist}
f
=
Dg
&=
Da_0(\vartheta)
+
\sum_{k=1}^m
\sum_{i=1}^{r_k}
D
\frac{a_{ik}(\vartheta)}{b_k(\vartheta)^i}.
\label{buch:intalg:liouville:speziell:eqn:fabgeleitet}
\intertext{Dabei ist wichtig, dass $f\in F$ ist und daher die Erweiterung
$\vartheta$ nicht enthält.
Wir wollen zeigen, dass diese zur Folge hat, dass gar keine Bruchterme
vorkommen können.
Die Ableitung eines einzelnen Terms in der Summe ist}
D
\frac{a_{ik}(\vartheta)}{b_k(\vartheta)^i}
&=
\frac{
Da_{ik}(\vartheta)\cdot b_k(\vartheta)^i
-
a_{ik}(\vartheta)ib_k(\vartheta)^{i-1} Db_k(\vartheta)
}{
b_k(\vartheta)^{2i}
}
\notag
\\
&=
\frac{Da_{ik}(\vartheta)}{b_k(\vartheta)^i}
-
i
\frac{
a_{ik}(\vartheta) Db_k(\vartheta)
}{
b_k(\vartheta)^{i+1}
}.
\label{buch:intalg:liouville:speziell:eqn:abgeleitet}
\end{align}
Die Ableitung $Db_k(\vartheta)$ hat Grad
$\deg Db_k(\vartheta) = \deg b_k(\vartheta)-1$.
Da $b_k(\vartheta)$ irreduzibel ist, sind $Db_k(\vartheta)$
und $b_k(\vartheta)$ teilerfremd.
Die Terme 
\eqref{buch:intalg:liouville:speziell:eqn:abgeleitet}
sind alle vollständig gekürzt.
Die höchste Nennerpotenz von $b_k(\vartheta)$ in
\eqref{buch:intalg:liouville:speziell:eqn:fabgeleitet}
hat den Exponenten $r_k+1$, und sie kommt nur genau einmal
im letzten Summanden vor.
Daher gibt es keinen weiteren Term, mit dem er sich wegheben könnte.
Die Summe in~\eqref{buch:intalg:liouville:speziell:eqn:fabgeleitet}
muss daher verschwinden und $g$ ist nur ein Polynom von $a_0(\vartheta)$.

Da $f=Dg=Da_0(\vartheta)$ nicht von $\vartheta$ abhängen darf, muss
$a_0(\vartheta)$ ein Polynom vom Grad höchstens $1$ sein, also
\[
g
=
a_0(\vartheta) = c_0+c_1\vartheta
\qquad\text{mit $c_k\in F$.}
\]
Damit $\vartheta$ in $Dg$ nicht vorkommt, muss der Grad von des Polynoms
$a_0(\vartheta)$ beim Differenzieren um eins abnehmen, was nur möglich ist,
wenn der Koeffizient von $\vartheta$ eine Konstante ist, also $c_1\in K$.
Somit 
\end{proof}

%
% Der Fall einer exponentiellen Erweiterung
%
\subsubsection{Der Fall einer exponentiellen Erweiterung}
Nach dem Prinzip von Liouville treten in einer Stammfunktion nur
logarithmische Erweiterungen auf.
Eine exponentielle Erweiterung $F(\vartheta)$ sollte also in der
Stammfunktion einer Funktion $f\in F$ gar nicht vorkommen.
Dies zeigt auch der folgende Satz.

\begin{satz}
\label{buch:intalg:liouville:satz:expsimple}
Sei $F(\vartheta)$ eine exponentielle transzendente Erweiterung eines
Funktionenkörpers $F$.
$F$ und $F(\vartheta)$ haben den gleichen Konstantenkörper.
Weiter sei $f\in F$ und $g\in F(\vartheta)$ eine Stammfunktion von $f$.
Dann ist $g\in F$.
\end{satz}

\begin{proof}
Sei also $\vartheta = \exp(u)$ eine Exponentialfunktion von $u\in F$.
Die Stammfunktion $g\in F(\vartheta)$ kann wieder mit einer Zerlegung
in irreduzible Faktoren des Nenners in der Form einer Partialbruchzerlegung
\eqref{buch:intalg:liouville:speziell:eqn:f}
geschrieben werden.

Im Fall einer logarithmischen Erweiterung
(Beweis von Satz~\ref{buch:intalg:liouville:satz:logsimple})
haben wir dann unter Verwendung von
Satz~\ref{buch:intalg:liouville:polynome:satz:logarithmisch}
gezeigt, dass der Term mit der höchsten Nennerpotenz $b_k(\vartheta)^{r_k+1}$
sich gegen nichts wegheben kann.
Dies hat dann gezeigt, dass die Bruchterme in der Stammfunktion nicht
vorkommen können.

Im vorliegenden Fall ist $b_k(\vartheta)$ kein Monom, also nicht von
der Form $h\vartheta^m$.
Nach Aussage 3 von Satz~\ref{buch:intalg:liouville:polynome:satz:exponentiell}
sind also $b_k(\vartheta)$ und $b_k(\vartheta)'$ teilerfremd und damit
die Bruchterme in \eqref{buch:intalg:liouville:speziell:eqn:f}
vollständig gekürzt.
Wie im Fall einer logarithmischen Erweiterung folgt jetzt, dass sich
die Bruchterme nicht wegheben können.
Da auf der linken Seite von \eqref{buch:intalg:liouville:speziell:eqn:f}
$\vartheta$ nicht vorkommen darf, dürfen in der Stammfunktion keine 
Bruchterme vorkommen.

Es bleibt also nur noch der polynomielle Teil.
Nach Satz~\ref{buch:intalg:liouville:polynome:satz:exponentiell}
kann aber der Grad des Polynoms $a_0(\vartheta)$ beim Differenzieren
nicht abnehen, die Ableitung $Da_0(\vartheta)$ kann daher $\vartheta$
nicht enhalten, wenn $a_0(\vartheta)$ ein Polynom vom, Grad 0 ist,
also $g=a_0(\vartheta)=a_{0,0}\in F$.
\end{proof}

%
% Der Fall einer algebraischen Erweiterung
%
\subsubsection{Der Fall einer algebraischen Erweiterung}
Auch für eine algebraische Erweiterung besagt das Prinzip von Liouville,
dass die Stammfunktion einer Funktion des Funktionenkörpers die
Erweiterung gar nicht enthält.

\begin{satz}
\label{buch:intalg:liouville:satz:algsimple}
Sei $F(\vartheta)$ eine algebraische Erweiterung eines
Funktionenkörpers $F$ mit Minimalpolynom $m(X)\in F[X]$.
$F$ und $F(\vartheta)$ haben den gleichen Konstantenkörper.
Weiter sei $f\in F$ und $g\in F(\vartheta)$ eine Stammfunktion von $f$.
Dann ist $g\in F$.
\end{satz}

\begin{proof}
Sei jetzt $\vartheta$ ein algebraisches Element über $F$ mit Minimalpolynom
$p$ mit Koeffizienten in $F$.
Wie früher ist die Stammfunktion $g$ eine rationale Funktion 
in $\vartheta$ mit Koeffizienten in $F$.
Nach Satz~\ref{buch:intalg:liouville:polynome:satz:algebraisch}
lässt sie sich aber auch als Polynom in $\vartheta$ mit Koeffizienten
in $F$ schreiben, wir schreiben $g=a(\vartheta)$ mit $u(Z)\in F[Z]$.
Die Ableitung ist
\[
f
=
Dg
=
Da(\vartheta).
\]
Die Ableitung des Polynoms $a$ in $\vartheta$ mit Koeffizienten
in $F$ ist
\begin{align*}
Da(\vartheta)
&=
\sum_{k=0}^m Da_k\vartheta^k
+
\sum_{k=1}^m ka_k\vartheta^{k-1} D\vartheta
\\
&=
\sum_{k=0}^{m-1}
\bigl(
Da_k
+
(k+1)a_{k+1}
D\vartheta
\bigr)
\vartheta^k
+
Da_m\vartheta^m.
\end{align*}
Darin ist $D\vartheta$ die Ableitung von $\vartheta$, die nach
Satz~\ref{buch:intalg:liouville:polynome:satz:algebraisch}
eine rationale Funktion von $\vartheta$ mit Koeffzienten in $F$ ist.
Es folgt, dass $Dg$ eine rationale Funktion von $\vartheta$ mit
Koeffizienten in $F$ ist, die ihrerseits wieder als Polynom
geschrieben werden kann.

Es muss jetzt noch gezeigt werden, dass aus $f=Dg$ folgt,
dass $a(\vartheta)$ nicht von $\vartheta$ abhängt.
$\vartheta$ ist nicht die einzige Lösung von $m(\vartheta)=0$, wir schreiben
$\vartheta=\vartheta_1$ und $\vartheta_2,\dots,\vartheta_l$ für die
$l=\deg m$ weiteren Lösungen der Gleichung.
Es gilt also
\begin{align}
f
&=
Da(\vartheta_k),\qquad k=1,\dots,l.
\notag
\intertext{Die Summe all dieser Gleichungen ist}
m\cdot f
&=
\sum_{k=1}^l Da(\vartheta_k).
\notag
\intertext{Es folgt}
f
&=
\frac{1}{l}
\sum_{k=1}^l Da(\vartheta_k)
=
Dh
\qquad\text{mit } h=\frac{1}{l} \sum_{k=1}^l a(\vartheta_k)
\label{buch:intalg:liouville:alg:eqn:symmetrisch}
\end{align}
$h$ ist eine symmetrische Funktion in
$\vartheta_1,\dots,\vartheta_l$, nach dem Hauptsatz über symmetrische
Funktionen lassen sich die Koeffizienten von $h$ nach dem Fundamentalsatz
\ref{buch:komplex:fundamental:satz:fundamentalpolynom} über symmetrische
Polynome durch die elementarsymmetrischen Funktionen von
$\vartheta_1,\dots,\vartheta_l$ ausdrücken.
Diese sind jedoch Koeffzienten des Minimalpolynoms $m(X)$ von $\vartheta$,
die alle in $F$ liegen.
Somit hängt die rechte Seite von
\eqref{buch:intalg:liouville:alg:eqn:symmetrisch}
nicht von $\vartheta$ ab.
\end{proof}
